\section{Заключение}

После выполнения лабораторной работы у меня появились ответы на следующие вопросы.

\paragraph{Почему степень точности квадратурной формулы определяется через многочлены?} По аппроксимационной теореме Вейерштрасса, любую непрерывную на отрезке функцию можно сколь угодно точно приблизить алгебраическим многочленом. Следовательно, если квадратурная формула точно интегрирует многочлены высокой степени, то для гладких функций её погрешность будет определяться лишь погрешностью аппроксимации.

То есть степень точности описывает порядок малости погрешности метода при стремлении шага разбиения к нулю.

\paragraph{Какая асимптотическая сложность у алгоритма?} Как видно из блок-схемы, вычислительная сложность одной итерации метода трапеций \textit{линейная} относительно количества разбиений отрезка. Так как алгоритм последовательно удваивает мощность сетки, а погрешность метода убывает как $O(h^2)$, то общее число вычислений функции до достижения точности $\varepsilon$ оценивается как $\mathcal{O}(\varepsilon^{-1/2})$ \textit{(без учёта верхнего предела по мощности разбиения)}.

Пространственная сложность алгоритма в идеале равна $\mathcal{O}(1)$: в памяти могут храниться лишь переменные для аккумулирования суммы, границы отрезков и текущий шаг. Остальные структуры данных, которые используются в реализации, нужны преимущественно для тестирования алгоритма.

\paragraph{В чём преимущество метода трапеций?} Главное преимущество метода в его простоте: для вычисления интеграла достаточно лишь последовательно суммировать значения функции в узлах сетки. Методу не нужно вычислять производные или решать дополнительные системы. Ещё можно добавить, что метод трапеций точно интегрирует линейные функции при одном разбиении.

\paragraph{Какие недостатки обнаружены у метода трапеций?} Основной недостаток --- низкий порядок аппроксимации ($O(h^2)$). Для достижения высокой точности требуется серьёзное увеличение числа разбиений, что ведёт к росту вычислительных затрат и накоплению ошибок округления.

Метод также неэффективен для функций, у которых производная быстро меняется: кусочно-линейная аппроксимация плохо описывает такое поведение и требует сетки большей мощности.

Если сравнивать его с методом Симпсона, метод трапеций требует больше вычислений функции для достижения сопоставимой точности на гладких участках \textit{(у последнего метода алгебраическая степень точности --- 3)}.

\paragraph{Почему количество разбиений для функций с особенностями существенно возрастает?} Метод трапеций основан на кусочно-линейной аппроксимации подынтегральной функции. В окрестности точки разрыва первого рода или устранимой особенности поведение функции резко отличается от линейного, что приводит к локальному росту погрешности аппроксимации.

Для выполнения глобального критерия останова $|I_k - I_{k-1}| < \varepsilon$ алгоритму нужно многократно удваивать число разбиений, пока шаг $h$ не станет достаточно малым, чтобы вклад ошибки в окрестности особенности уменьшился ниже заданного порога.